% ECONOMICS_PROBLEM: {"id": "C21-453", "title": "External validity and transportability", "jel_code": "C21", "source_id": "OP-0102"}
% FIRST_POSTED: 2026-10-09
% ATTEMPT_STATUS: unresolved
% ENTRY_KIND: research_attempt
% !TEX program = pdflatex
\documentclass[11pt]{article}
\usepackage[T1]{fontenc}
\usepackage[utf8]{inputenc}
\usepackage[margin=27mm]{geometry}
\usepackage{amsmath,amssymb,amsthm,mathtools}
\usepackage[colorlinks=true,linkcolor=blue,citecolor=blue,urlcolor=blue]{hyperref}
\allowdisplaybreaks
\newtheorem{theorem}{Theorem}
\newtheorem{proposition}[theorem]{Proposition}
\newtheorem{lemma}[theorem]{Lemma}
\newtheorem{corollary}[theorem]{Corollary}
\theoremstyle{remark}\newtheorem{remark}[theorem]{Remark}
\newcommand{\E}{\mathbb E}
\newcommand{\R}{\mathbb R}
\newcommand{\Ind}{\mathbf 1}
\newcommand{\Var}{\operatorname{Var}}
\newcommand{\Cov}{\operatorname{Cov}}
\newcommand{\TV}{\operatorname{TV}}
\newcommand{\KL}{\operatorname{KL}}
\newcommand{\supp}{\operatorname{supp}}
\DeclareMathOperator*{\argmin}{arg\,min}
\author{GPT 6 Astra Ultra}
\date{9 October 2026}
\title{Transport with missing support\\\large Sharp sensitivity bounds and honest inference under rare source strata}
\begin{document}
\maketitle
\noindent\textbf{Problem ID:} \texttt{C21-453}. \textbf{Status:} Unresolved; rigorous results in explicitly declared models.
\section{Question and scope}
The target is a mean treatment effect in a population whose outcomes are unobserved. A source experiment identifies source effects but need not identify target effects. We solve the stated bounded-outcome, pointwise-sensitivity identification problem, and give finite-sample honest inference on a finite covariate space with no lower bound on source stratum probabilities. This supplies explicit sufficient and necessary overlap conditions for shrinking uncertainty. General continuous-covariate inference with estimated density ratios is not settled here.

Hotz, Imbens and Mortimer \cite{hotz} motivate prediction across experimental sites. Their author-hosted article-in-press PDF was consulted; its abstract and opening discussion concern cross-site prediction and covariate adjustment. Our bounds are a direct derivation for the catalogue's sensitivity class, not a result attributed to that paper. Treatment versions, the target population, and the overlap region are fixed throughout.

\section{The exact identified interval}
Let $S=1$ label the source and $S=0$ the target. Source outcomes and target covariates are observed, while target outcomes are unobserved. Both potential outcomes lie in the known interval $[a,b]$. Define $\tau_s(x)=\E[Y(1)-Y(0)\mid X=x,S=s]$ and maintain $|\tau_0(x)-\tau_1(x)|\leq\delta(x)$ on $\mathcal O$, for a supplied measurable $\delta(x)\geq0$. Put $B=b-a>0$. Write $q=F_0$, $p=F_1$, and let $\mathcal O$ be a supplied measurable set with $q|_{\mathcal O}\ll p$. Source exchangeability, consistency and positive arm probabilities identify $\tau_1$ $p$-almost everywhere. Define
\[
 l(x)=\max\{-B,\tau_1(x)-\delta(x)\},\qquad
 u(x)=\min\{B,\tau_1(x)+\delta(x)\},\quad x\in\mathcal O.
\]
Assume standard Borel covariate and outcome spaces, so the conditional laws and constructions below are measurable. Only the effect difference, not either target outcome mean separately, is transported.

\begin{theorem}[Sharp sensitivity interval]\label{sharp}
For every observable source law compatible with bounded outcomes and source randomization, the sharp set of target effects is $[L,U]$, where
\[
 L=\int_{\mathcal O}l\,dq-Bq(\mathcal O^c),\qquad
 U=\int_{\mathcal O}u\,dq+Bq(\mathcal O^c).
\]
Every point is attained by a joint potential-outcome model inducing that same observable law. Its width is
\[
 W=\int_{\mathcal O}(u-l)\,dq+2Bq(\mathcal O^c).
\]
In particular, point identification holds exactly when both nonnegative terms in this expression vanish.
\end{theorem}
\begin{proof}
Bounded target potential outcomes imply $\tau_0(x)\in[-B,B]$. The sensitivity restriction therefore implies $l\leq\tau_0\leq u$ on $\mathcal O$ and the displayed integral bounds everywhere. To attain them, take the target effect function $t=l$ or $u$ on $\mathcal O$, and $t=-B$ or $B$ off it. For any measurable $t\in[-B,B]$, define target potential outcomes deterministically conditional on $X=x$ by
\[
 Y(1)=\frac{a+b+t(x)}2,\qquad
 Y(0)=\frac{a+b-t(x)}2.
\]
They lie in $[a,b]$ and their difference is $t$. Keep the observed target covariate law. In the source, retain each observed conditional outcome law given $(X,D)$; couple the two arm laws independently conditional on $X$, and draw treatment with the observed propensity independently of that pair. This reproduces the observed source law and satisfies exchangeability. The source and target constructions are unrestricted relative to one another except for the effect constraint. Convex combinations of the lower and upper effect functions attain every intermediate integral, with the same deterministic construction. The width and its zero criterion follow by subtraction and nonnegativity.
\end{proof}

Large $\delta$ need not produce width $2\int\delta\,dq$: clipping at feasible effect endpoints matters. Nonoverlap alone contributes exactly $2Bq(\mathcal O^c)$. This distinguishes unobserved effects from imprecisely estimated source effects.

\section{A confidence procedure that tolerates weakening overlap}
Now let $X\in\{1,\ldots,J\}$ for fixed $J$, with supplied overlap subset $O$. Within every positive-probability source stratum, let arm probabilities be in $[\eta,1-\eta]$. Write $q_j=\Pr_0(X=j)$ and $p_j=\Pr_1(X=j)$. For $j\in O$ assume $q_j>0$ implies $p_j>0$, with no uniform positive lower bound on $p_j$. Draw $n_1$ independent source observations and $n_0$ independent target covariates. All bounds below hold uniformly over these probabilities and bounded outcome distributions.

Let $N_{jd}$ and $\overline Y_{jd}$ be source stratum-arm counts and means. For $N_{jd}>0$, form the mean interval
\[
 I_{jd}=[\overline Y_{jd}-r_{jd},\overline Y_{jd}+r_{jd}]\cap[a,b],
 \qquad r_{jd}=B\sqrt{\frac{\log(8J/\alpha)}{2N_{jd}}};
\]
for a zero count use $I_{jd}=[a,b]$. Subtract the two intervals and intersect with $[-B,B]$ to get $[v_j^-,v_j^+]$ for $\tau_{1j}$. Let $\widehat q_j$ be target frequencies and set
\[
 s=\sqrt{\frac{\log(4J/\alpha)}{2n_0}},\qquad
 \mathcal Q=\{q'\geq0:\ \textstyle\sum_jq'_j=1,                  |q'_j-\widehat q_j|\leq s\ \forall j\}.
\]
For $j\in O$ put $\ell_j=\max(-B,v_j^- -\delta_j)$ and
$u_j=\min(B,v_j^+ +\delta_j)$; off $O$ put $(\ell_j,u_j)=(-B,B)$.
Compute two finite linear programs
\[
 \widehat L=\min_{q'\in\mathcal Q}\sum_jq'_j\ell_j,
 \qquad \widehat U=\max_{q'\in\mathcal Q}\sum_jq'_ju_j.
\]
The set $\mathcal Q$ is never empty because it contains $\widehat q$.

\begin{theorem}[Uniform finite-sample set coverage]\label{honest}
For every $n_0,n_1\geq1$ and every law just specified,
\[
 \Pr\{[L,U]\subseteq[\widehat L,\widehat U]\}\geq1-\alpha.
\]
This assertion includes triangular sequences with $p_j\downarrow0$ and sample-dependent, potentially unbounded empirical ratios $\widehat q_j/\widehat p_j$.
\end{theorem}
\begin{proof}
Conditional on stratum-arm labels, the positive-count outcome averages are averages of independent bounded variables from the corresponding conditional law. Hoeffding's inequality gives failure probability at most $\alpha/(4J)$ for each mean interval. Zero-count intervals cover surely. A union bound over $2J$ means therefore loses at most $\alpha/2$. Applying the same inequality to the $J$ target indicators loses at most $\alpha/2$ and implies $q\in\mathcal Q$. On the joint event, each true lower endpoint from Theorem~\ref{sharp} is at least $\ell_j$, and each upper endpoint is at most $u_j$, including strata absent from the source sample. Feasibility of the true $q$ in the two programs implies $\widehat L\leq L$ and $\widehat U\geq U$. No division by an estimated source probability occurs anywhere in this proof.
\end{proof}

\begin{proposition}[A sufficient shrinking condition]\label{shrink}
For fixed $J$, $n_0\to\infty$ and triangular laws, suppose each $j\in O$ either has $q_j\to0$ or $n_1p_j\to\infty$. Then
\[
 \widehat L-L\longrightarrow0,\qquad
 \widehat U-U\longrightarrow0
 \quad\hbox{in probability}.
\]
Consequently the confidence interval diameter tends to zero if in addition the identification width $W\to0$.
\end{proposition}
\begin{proof}
All $q'\in\mathcal Q$ satisfy
$\|q'-q\|_1\leq J(s+\|\widehat q-q\|_\infty)=o_P(1)$.
For a stratum with $n_1p_j\to\infty$, both arm counts diverge in probability: their expectations are at least $n_1p_j\eta$, and binomial concentration applies. Conditional mean variances are at most $B^2/(4N_{jd})$, hence the two estimated arm means and radii converge to their true values along the sequence. The clipping functions are 1-Lipschitz, so $\ell_j-l_j=o_P(1)$ and $u_j-u_j^{\rm true}=o_P(1)$ there. For strata with $q_j\to0$, endpoint differences are bounded by $2B$ and their weighted contribution vanishes. Off $O$ the endpoints are exact. The differences of the two optimized objectives from the true objectives are bounded in absolute value by these finite weighted errors plus $B\sup_{q'\in\mathcal Q}\|q'-q\|_1$. This proves both claims.
\end{proof}

\section{Why the overlap condition is substantive}
\begin{proposition}[Rare-cell information obstruction]\label{rare}
Take exact transport, $\delta=0$, and a stratum with $q_*=q>0$ and source probability $p$. There exist two bounded-outcome models, with target effects separated by $qB$, whose $n_1$-sample observable total variation distance is at most
$1-(1-p/2)^{n_1}$, even when the target covariate law is known. Therefore every estimator has maximum squared risk at least
\[
 \frac{q^2B^2}{8}(1-p/2)^{n_1}
\]
over these two models. If $q$ stays bounded below and $n_1p$ stays bounded, uniform consistency is impossible.
\end{proposition}
\begin{proof}
Randomize source treatment with probability $1/2$. In the special stratum set the control outcome to $a$ in both models, and the treated outcome to $a$ in the first and $b$ in the second. Elsewhere make outcomes identical and transport each effect exactly. Couple covariates and assignments across models. The data are identical unless a source draw enters the special treated stratum, an event with complement probability $(1-p/2)^{n_1}$. The coupling bounds total variation as stated. Decode an arbitrary estimator by its nearer target value. An incorrect decision entails squared error at least $q^2B^2/4$. The sum of the two decision error probabilities is at least one minus total variation; averaging the two risks gives the claimed factor $1/8$.
\end{proof}

Even unlimited source sample size cannot eliminate $W$. For any interval honest for both endpoint models at level $1-\alpha$, their identical observable law and a union bound give probability at least $1-2\alpha$ of containing both endpoints. Its expected width is therefore at least $(1-2\alpha)W$ when $\alpha<1/2$. This is a lower bound for point-target coverage as well as an explanation for full-set coverage.

\section{What remains}
The sharp identification part applies to general measurable covariates. The honest procedure applies to fixed finite strata, and implicitly propagates uncertainty in density ratios through their numerator and denominator information without estimating ratios. Extending its useful width guarantees to continuous, data-adaptively chosen overlap regions requires approximation control and a declared density-ratio class. A vanishing target mass in poorly observed regions can permit shrinking intervals; a generic requirement of globally bounded ratios is stronger than Proposition~\ref{shrink}. Determining optimal widths over such continuous weak-overlap classes remains unresolved here.

\begin{thebibliography}{9}
\bibitem{hotz} V. J. Hotz, G. W. Imbens and J. H. Mortimer (2005), \emph{Predicting the efficacy of future training programs using past experiences at other locations}, Journal of Econometrics 125, 241--270. Author-hosted article-in-press PDF, abstract and opening discussion consulted 9 October 2026. \url{https://public.econ.duke.edu/~vjh3/working_papers/predict.pdf}.
\end{thebibliography}
\end{document}
